\documentclass{report}
\usepackage{setspace}
%\usepackage{subfigure}

\pagestyle{plain}
\usepackage{amssymb,graphicx,color}
\usepackage{amsfonts}
\usepackage{latexsym}
\usepackage{a4wide}
\usepackage{amsmath}

\newtheorem{theorem}{THEOREM}
\newtheorem{lemma}[theorem]{LEMMA}
\newtheorem{corollary}[theorem]{COROLLARY}
\newtheorem{proposition}[theorem]{PROPOSITION}
\newtheorem{remark}[theorem]{REMARK}
\newtheorem{definition}[theorem]{DEFINITION}
\newtheorem{fact}[theorem]{FACT}

\newtheorem{problem}[theorem]{PROBLEM}
\newtheorem{exercise}[theorem]{EXERCISE}
\def \set#1{\{#1\} }

\newenvironment{proof}{
PROOF:
\begin{quotation}}{
$\Box$ \end{quotation}}



\newcommand{\nats}{\mbox{\( \mathbb N \)}}
\newcommand{\rat}{\mbox{\(\mathbb Q\)}}
\newcommand{\rats}{\mbox{\(\mathbb Q\)}}
\newcommand{\reals}{\mbox{\(\mathbb R\)}}
\newcommand{\ints}{\mbox{\(\mathbb Z\)}}

%%%%%%%%%%%%%%%%%%%%%%%%%%


\title{  	{ \includegraphics[scale=.5]{ucl_logo.png}}\\
{{\Huge HYPERDAPO}}\\
{\large A DAPO-based framework for optimal execution of Hyperliquid limit orders}\\
		}
\date{Submission date: 08 September 2025}
\author{Mohammed Khalil\thanks{
{\bf Disclaimer:}
This report is submitted as part requirement for the MSc Information Security at UCL. It is
substantially the result of my own work except where explicitly indicated in the text.
\newline  %% \\ screws it up
The report will be distributed to the internal and external examiners, but thereafter may not be copied or distributed except with permission from the author.}
\\ \\
MSc Information Security\\ \\
Arthur Gervais}



\begin{document}
 
\onehalfspacing
\maketitle
\begin{abstract}
Summarise your report concisely.
\end{abstract}
\tableofcontents
\setcounter{page}{1}


\chapter{title of first chapter}
This is just a bare minimum to get started.  There is unlimited guidance on using latex, e.g. {\tt https://en.wikibooks.org/wiki/LaTeX}.   You are still responsible to check the detailed requirements of a project, including formatting instructions, see \\
{\tt https://moodle.ucl.ac.uk/pluginfile.php/3591429/mod\_resource/content/7/UGProjects2017.pdf}.
Leave at least a line of white space when you want to start a new paragraph.

Mathematical expressions are placed inline between dollar signs, e.g. $\sqrt 2, \sum_{i=0}^nf(i)$, or in display mode
\[ e^{i\pi}=-1\] and another way, this time with labels,
\begin{align}
\label{line1} A=B\wedge B=C&\rightarrow A=C\\
&\rightarrow C=A\\
\intertext{note that}
n!&=\prod_{1\leq i\leq n}i \\
\int_{x=1}^y \frac 1 x \mathrm{d}x&=\log y
\end{align}
We can refer to labels like this \eqref{line1}.   

\chapter{title of second chapter}
Often lots of citations here (and elsewhere), e.g. \cite{Rey:D} or \cite[Theorem 2.3]{PriorNOP70}.   Bibtex can help with this, but is not essential. If you want pictures, try

You can use 
\begin{itemize}
\item lists
\item like this
\end{itemize}
or numbered
\begin{enumerate}
\item like this,
\item or this
\end{enumerate}
but don't overdo it.
\chapter{title of third chapter}
If you have a formal theorem you might try this.
\begin{definition}\label{def}
See definition~\ref{def}.
\end{definition}
\begin{theorem}
For all $n\in\nats,\; 1^n=1$.
\end{theorem}
\begin{proof}
By induction over $n$.
\end{proof}

\chapter{etc.}
\appendix


\begin{thebibliography}{HHM99}


\bibitem[Pri70]{PriorNOP70}  %%only an example
A.~Prior.
\newblock The notion of the present.
\newblock {\em Studium Generale}, 23:  245--248, 1970.


\bibitem[Rey97]{Rey:D}
M.~Reynolds.
\newblock A decidable temporal logic of parallelism.
\newblock {\em Notre Dame Journal of Formal Logic}, 38(3):  419--436,
  1997.
\end{thebibliography}
\chapter{Other appendices, e.g. code listing}

\end{document}